\documentclass{article}
\usepackage[T1]{fontenc}
\usepackage{lmodern}
\usepackage{graphicx}
\usepackage{amsmath,amssymb}
\usepackage[a4paper,margin=2cm]{geometry}
\usepackage[absolute,overlay]{textpos}
\setlength{\TPHorizModule}{1mm}
\setlength{\TPVertModule}{1mm}
\usepackage[indonesian]{babel}
\usepackage{hyperref}
\setlength{\emergencystretch}{2em}
% unit: Halaman 5

\begin{document}

% === Halaman 5 (llm) ===

\section*{Apa Steganografi itu?}

\begin{itemize}
    \item Dari Bahasa Yunani: \textit{steganos} + \textit{graphien}
    \begin{itemize}
        \item \textcolor{red}{"steganos"} ($\sigma\tau\epsilon\gamma\alpha\nu\acute{o}\varsigma$): tersembunyi
        \item \textcolor{red}{"graphien"} ($\gamma\rho\alpha\phi\acute{\iota}\alpha$): tulisan
    \end{itemize}
    steganografi: tulisan tersembunyi (\textit{covered writing})
    
    \item \textbf{Steganography}: ilmu dan seni menyembunyikan pesan rahasia dengan suatu cara sedemikian sehingga tidak seorang pun yang mencurigai keberadaan pesan tersebut.
\end{itemize}

Tujuan steganografi: pesan tidak terdeteksi keberadaannya

\end{document}
